\section{Introduction}
\label{sec:introduction}
High-dimensional survival analysis seeks interpretable associations between incomplete event times and many dependent covariates. Molecular measurements often combine a low-dimensional common component with feature-specific variation. In an accelerated failure time (AFT) model, the scientific target is a sparse slope on the observed covariates. Factor adjustment offers a representation of their dependence, but the factors used in estimation are themselves learned from the sample. A statistical analysis must therefore connect the original slope to the estimated factor coordinates while accounting for right censoring. We develop that connection for a smoothed rank estimator.

Rank-based AFT regression estimates conditional log-time associations with an unspecified error distribution \citep{Wei1992,JinEtAl2003}. Regularized Gehan estimation, joint clinical adjustment, and scalable high-dimensional algorithms are developed by \citet{CaiEtAl2009}, \citet{Johnson2009}, and \citet{SuderMolstad2022}. Distribution-function and parameter-perturbation smoothing produce differentiable censored-rank criteria \citep{Heller2007,JohnsonStrawderman2009}. Related computational and robust formulations include \citet{ChoiChoi2021}, \citet{LiEtAl2021}, and Seon et~al.\ (2026, preprint, arXiv:2606.31094). Our estimator uses a common scalar Gaussian convolution of the positive-part loss, so the bandwidth enters the analysis through an explicit smoothing-bias term.

Factor adjustment retains the common component as an unpenalized regression block. FarmSelect develops this construction for regularized estimation \citep{FanEtAl2020}; Bayle and Fan (2022, preprint, arXiv:2210.01067) study factor-augmented Cox regression. For censored rank estimation, the estimated factor representation and the patient pairs in the loss share the same covariate sample. This dependence requires a joint treatment of score centering, pairwise fluctuations, smoothing bias, and estimated-design curvature. The unpenalized clinical and factor coefficients must also be included in the error geometry.

Our primary target is $\bm{\beta}_0$ in the ordinary observed-covariate model
\begin{equation}
 \begin{aligned}
 \bm{X}_i&=\bm{B}\bm{f}_i+\bm{u}_i,\\
 \log T_i&=\bm{Z}_i^\top\bm{\theta}_0+\bm{X}_i^\top\bm{\beta}_0+\varepsilon_i,
 \end{aligned}
 \label{eq:intro-model}
\end{equation}
where $\bm{X}_i\in\mathbb R^p$, $\bm{Z}_i\in\mathbb R^d$, and the factor dimension $K$ is small. The predictor is equivalently $\bm{Z}_i^\top\bm{\theta}_0+\bm{u}_i^\top\bm{\beta}_0+\bm{f}_i^\top\bm{\gamma}_0$, with $\bm{\gamma}_0=\bm{B}^\top\bm{\beta}_0$. We estimate the common and residual components by principal component analysis (PCA), then minimize a smoothed Gehan loss with an $\ell_1$ penalty on $\bm{\beta}$. Clinical and factor coefficients are estimated jointly and remain unpenalized. The estimator allows the factor coefficient to vary freely; it contains the ordinary fitted predictor as a special case while retaining $\bm{\beta}_0$ as its sparse target. Section~\ref{sec:model} also considers an additional factor effect relative to a specified latent representation.

The main contribution is a finite-sample analysis with same-sample factor estimation. An exact reference representation preserves the centered ordinary-model predictor and separates estimated-design perturbation from the approximation associated with an additional factor effect. Conditional on the full design, an event-weighted residual-difference identity centers the unsmoothed score. Pairing and contraction bound the empirical contrasts at the design's column root-mean-square scale, avoiding the extra logarithmic loss of an entrywise-envelope argument. Local event information and a regularization argument that includes the unpenalized block yield coefficient bounds and finite attainment on the same statistical event. Primitive factor conditions supply the required estimated-design geometry. The analysis uses integrable loss differences and accommodates log-time errors without moments under the stated density and censoring conditions.

A second contribution connects statistical and numerical accuracy. A complete stationarity residual transfers the coefficient guarantee to finite numerical candidates on the same fixed-tuning event. Event-row blocking evaluates the full pairwise gradient efficiently. We provide a reference proximal-gradient convergence result and implement the same objective using accelerated iterations with qualified active-face Newton refinement. These components separate the statistical target, the search algorithm, and the accuracy required of its output.

The numerical illustrations connect this framework to computational geometry and training selection. In three fixed strong-dependence samples, factor adjustment has lower maximum active curvature and reaches common stationarity thresholds with fewer gradient evaluations. Residual projections identify weak-curvature directions in the late iterates; one Newton step completes each of the three unadjusted fits on its existing active face. A shared fold-specific selection procedure produces twelve full-training refits at the prescribed tolerance. The selected models also distinguish sparse residual effects from the dense molecular slope induced by the common factors. Each illustration addresses the numerical or representational role of the estimator; independent repeated-sample risk and biomedical applications are separate evaluation targets.

Section~\ref{sec:model} defines the observation scheme and regression targets. Section~\ref{sec:estimation} constructs the estimator and frozen prediction map. Sections~\ref{sec:theory} and~\ref{sec:computation} develop the statistical and numerical results. Section~\ref{sec:numerical} presents the computational and training-selection illustrations, and Section~\ref{sec:discussion} discusses their implications. Online Resource~1 contains the proofs in Supplementary Appendices~\ref{app:identification}--\ref{app:computation} and the reproducibility and implementation details in Appendices~\ref{app:reproducibility}--\ref{app:implementation}.
